% Player handouts, sideways maps and the handout appendix (#5)
\documentclass[letterpaper,twoside,twocolumn,openany,nodeprecatedcode]{dndbook}
\begin{document}
\tableofcontents
\DndListOfHandouts

\chapter{Handouts}

A letter shown here and again, full page, in the appendix as \DndHandoutRef{mayor}:

\begin{DndHandout}[style=letter, title={The Mayor's Letter}, label=mayor, appendix]
  To whoever finds this,

  The bell rang again last night. I fear the chapel has woken. Come quickly, and bring silver.

  \hfill --- Mayor Oswin
\end{DndHandout}

A torn note, inline only:

\begin{DndHandout}[style=note, title={Scrap of Paper}, label=scrap]
  Third pew from the altar. Under the stone.
\end{DndHandout}

A poster, only in the appendix (\DndHandoutRef{wanted}):

\begin{DndHandout}[style=poster, title=Wanted, label=wanted, art=examples/art/chapter, inline=false, appendix]
  \textbf{Grell the Butcher}\\
  For crimes against the crown.\\
  Reward: 500 gold pieces.
\end{DndHandout}

\begin{DndMap}[caption=The Ruins, label=ruins]{img/example-map}
  \DndMapArea{0.225,0.27}{Gatehouse}
\end{DndMap}

\DndArea{Gatehouse}
Text.

% The player version goes to the appendix as a handout
\DndPlayerMap[appendix]{ruins}

The player map is \DndHandoutRef{ruins}.

% A landscape map on its own page, turned sideways
\begin{DndMap}[caption=The Coast, label=coast, placement=sideways]{examples/art/landscape}
\end{DndMap}

Text after the sideways map.

\appendix
% Starts with \chapter{Handouts}; the first handout shares its page
\DndHandoutAppendix

\end{document}
